\section{Reparto del trabajo}

El desarrollo se ha organizado de forma que cada integrante asumiera la
responsabilidad principal de un conjunto de áreas del proyecto. La
Tabla~\ref{tab:reparto} lo resume y los apartados siguientes lo detallan.

\begin{table}[h]
\centering
\begin{tabular}{lp{8.6cm}}
\toprule
\textbf{Integrante} & \textbf{Áreas principales} \\
\midrule
Pablo Martín Palomino
    & Base del agente y modelo local, motor visual, analítica avanzada,
      evaluación y robustez del flujo de Bizum \\[0.2cm]
Natalia Yue Gallardo Pretel
    & Flujo conversacional de Bizum, seguridad y experiencia en la interfaz,
      y redacción de la memoria \\[0.2cm]
Pablo Reyes Sousa
    & Interacción por voz y personalización de la voz del asistente \\
\bottomrule
\end{tabular}
\caption{Resumen del reparto del trabajo.}
\label{tab:reparto}
\end{table}

\subsection{Pablo Martín Palomino}

\begin{itemize}
    \item Base inicial del proyecto: servidor FastAPI con WebSocket, ejecución
    de consultas SQL en modo solo lectura, primeros gráficos con Vega-Lite y
    voz en el navegador.
    \item Planificación del trabajo y migración a un modelo local
    (\texttt{qwen3:8b} sobre Ollama), incluido el script de arranque que amplía
    el contexto a 8192 tokens.
    \item Detección de pagos recurrentes y control del tamaño del contexto del
    modelo.
    \item Proyección del gasto a fin de mes, con el margen de error de cada
    categoría medido sobre el histórico.
    \item Motor visual: llamada dedicada para generar la especificación,
    decisión determinista en el backend sobre cuándo mostrar un elemento visual
    y tablas como segunda representación.
    \item Banco de evaluación de 44 preguntas y pruebas automáticas del flujo
    de Bizum y del detector de recurrencia.
    \item Robustez del flujo de Bizum: validación del importe antes de pedir el
    PIN, unificación de las dos vías de entrada, transacción
    \texttt{BEGIN IMMEDIATE} y acceso a la base de datos fuera del bucle de
    eventos.
    \item Reducción de la latencia: \textit{streaming} por frases, recorte del
    \textit{system prompt} y generación visual posterior a la respuesta hablada.
    \item Documentación del repositorio en español e inglés.
\end{itemize}

\subsection{Natalia Yue Gallardo Pretel}

\begin{itemize}
    \item Atajo para consultar el saldo sin llamar al modelo y reducción del
    número de tokens y llamadas por respuesta.
    \item Flujo conversacional de Bizum: validación del contacto, herramienta de
    contactos, importe opcional y normalización de las respuestas afirmativas y
    negativas en la confirmación.
    \item Primera versión del historial de conversación y corrección de las
    consultas de gasto, incluida la eliminación del formato Markdown en las
    respuestas.
    \item Capa de seguridad en la interfaz: pantalla de acceso, teclado numérico
    para el PIN y aviso de los intentos restantes.
    \item Experiencia de usuario: justificante de Bizum con concepto,
    personalización del perfil, del color y del tamaño de letra, y logotipos de
    los comercios en las tablas.
    \item Estructura de la memoria en \LaTeX{}, portada y redacción de los
    capítulos.
\end{itemize}

\subsection{Pablo Reyes Sousa}

\begin{itemize}
    \item Mejora de la interacción por voz: activación del micrófono con la
    barra espaciadora e indicador visual de que la lectura en voz alta está
    activada.
    \item Panel de personalización de la voz del asistente: elección entre las
    voces disponibles en español, ajuste del tono y de la velocidad, y prueba de
    la configuración antes de guardarla.
\end{itemize}

\subsection{Uso de asistentes de IA}

Durante el desarrollo se han utilizado asistentes de programación basados en
inteligencia artificial como apoyo en la escritura y revisión de código, la
depuración de errores, la redacción de documentación y la generación del
catálogo de comercios de los datos ficticios. Su uso ha estado siempre sujeto a la revisión del equipo.
